\documentclass[11pt]{article}
\usepackage[T1]{fontenc}
\usepackage{lmodern,amsmath,amssymb,microtype}
\usepackage[margin=1.02in]{geometry}
\usepackage[hidelinks]{hyperref}
\urlstyle{same}
\setlength{\emergencystretch}{2em}
\title{Which fields have algebraic closures in the finite partition model?}
\author{Wentao Li\\{\small Prepared with assistance from OpenAI Codex}}
\date{November 19, 2026 \quad $\cdot$ \quad Research draft}
\begin{document}
\maketitle
\section*{Question and result}
Amitayu Banerjee asked whether algebraic closures, when they exist, are unique up to isomorphism in Bruce's finite partition model \cite{ban}. This is Form 233 in the Howard--Rubin catalogue. In this model the ambient group is the full symmetric group on countably infinitely many atoms. A finite partition $P$ supports an object if $G_P$, the group of permutations preserving each block setwise, fixes it. Hereditarily supported objects form the model.

\paragraph{Theorem.} Every algebraically closed field in Bruce's model has a well-orderable underlying set. An internal field has an algebraic closure exactly when it is well-orderable, and any two such closures are isomorphic over the base field. The proof works in any permutation model over an ambient universe with choice in which the basic support subgroups are generated by elements of square one.

\paragraph{The group and field ingredients.} Every permutation of a set in the ambient universe with choice is the product of two permutations of square one. On each cyclic orbit index its points by $\mathbb Z/m\mathbb Z$, or by $\mathbb Z$ on an infinite orbit, so the permutation is the shift $j\mapsto j+1$. The product of $j\mapsto-j$ followed by $j\mapsto1-j$ is that shift. Apply this on every orbit. If the permutation preserves each block of a partition, the two resulting factors do also.

We use the classical involution fact attributed to Pincus in Banerjee's Lemma 3.2: if $E$ is an algebraically closed field and $i^2=-1$ in $E$, any field automorphism of square one fixing $i$ is the identity. It follows as well from the Artin--Schreier theorem: a nontrivial involution has real closed fixed field and cannot fix a square root of $-1$. In positive characteristic a nontrivial involution on an algebraically closed field is impossible.

\paragraph{Proof.} Fix an internal algebraically closed field $E$ and choose one $i\in E$ with $i^2=-1$. Take a common support $P$ of its field structure and $i$. The field is algebraically closed externally: every polynomial over it is specified by a finite tuple of internal coefficients, and that tuple is in the model. Thus each square-one element of $G_P$ acts on $E$ by a field automorphism fixing $i$, which must be trivial by the involution fact. The generation property of $G_P$ gives pointwise fixation of $E$. Any ambient well-order of $E$ is consequently supported by $P$ and belongs to the model.

A field that embeds into such a field by an internal map inherits a well-order. Conversely, a well-order of a field $F$ gives an internal isomorphism with a field on an ordinal, by transferring the operations along the order isomorphism. Those operations are pure sets. The pure part of the permutation model is the pure part of the ambient universe and satisfies ZFC. Construct an algebraic closure of the ordinal copy there, and transport its embedding back to $F$. If the definition requires literal containment, relabel the copy of the base by $F$ and tag the other elements disjointly.

For uniqueness, let $j_1:F\to E_1$ and $j_2:F\to E_2$ be two internal algebraic closures. Being algebraic over the base is absolute too, by the same finite-coefficient argument. The ordinary uniqueness theorem in the ambient universe gives an isomorphism $\phi:E_1\to E_2$ satisfying $\phi j_1=j_2$. Combine supports fixing $E_1$ and $E_2$ pointwise; every ordered pair in the graph of $\phi$ is then fixed, so $\phi$ is internal. \hfill$\square$

The field lemma and the classical uniqueness theorem are prior. Banerjee applied the lemma to a different permutation model in Theorem 3.3. Here the model-specific step checks that after supporting one root of $-1$, the full block stabilizer is generated by involutions.
\begin{thebibliography}{9}
\bibitem{ban} A. Banerjee, \emph{Partition models, permutations of infinite sets without fixed points, and weak forms of AC}, 2021, Question 8.1 in \url{https://arxiv.org/abs/2109.05914v1}; see Lemma 3.2, Theorem 3.3 and Question 7.1 in \url{https://arxiv.org/abs/2109.05914v4}.
\bibitem{bruce} B. B. Bruce, \emph{A Permutation Model with Finite Partitions of the Set of Atoms as Supports}, Rose-Hulman Undergraduate Mathematics Journal 17 (2016). \url{https://scholar.rose-hulman.edu/rhumj/vol17/iss1/4/}.
\end{thebibliography}
\end{document}
